\section{Asleep to the Real World}

\begin{quotex}
God has described Himself as the Outwardly Manifest and the Inwardly Hidden.
\flright{\textsc{Ibn Arabi}, \emph{Bezels of Wisdom}}
\end{quotex}


There can be no difference between the symbol and the doctrine, since the symbol is based on doctrine and doctrine can be expressed in symbols. Often the symbolism erupts in the mind and only later is the doctrine understood. That is why chaos and the non-manifested are revealed in a poem\footnote{\url{https://gornahoor.net/?p=15803}} rather than in text.

Rene Guenon describes the Two Nights in the collection \emph{Initiation and Spiritual Realization}. The lower darkness is the Chaos, the state of indifferentiation that is the starting point of manifestation. This is the first night. Those confined in the affairs of the world, the dunya, can only see a reflection of the principle. To them, the non-manifested is the night, the second night. Their only contact is in the causal state of deep sleep, so they remain unaware of the true causes in the world.

On the other hand, for those aware of the unmanifest, it is the world that is in darkness. Its many appearances are but different modes of the sole Existence.

\paragraph{The Myth}


\textbf{Selene} is the Moon goddess; her name is derived from the Greek word for ``light''. She is described like this:

\begin{quotex}
From her immortal head a radiance is shown from heaven and embraces earth; and great is the beauty that ariseth from her shining light.
\end{quotex}


Selene was in love with the mortal Endymion. Although pulled by her love, Endymion was still attached to the world, as shown in his perpetual state of sleep. The Moon goddess Selene, in her Love for him, kisses him sweetly as he sleeps. Hence, only in his dreams does he escape the illusion of Chaos. This incites pleasant dreams in him, but the Queen of Night remains elusive as long as he sleeps.

\begin{quotex}
\begin{verse}
Divine Selene watched him from on high,\\
and slid from heaven to earth; for passionate love\\
drew down the immortal stainless Queen of Night.
\end{verse}
\flright{\textsc{Quintus Smyrnaeus}}
\end{quotex}


\paragraph{Plato}


\textbf{Plato} was the prophet of the unmanifest. Every object is the reflection of an idea. The object is a sense perceptible phenomenon but the idea is supersensible. Objects perish but idea are immortal. Therefore, objects are transient and illusory. Only the ideas are real.

But the ideas express a unity. This higher unity, the fundamental idea, is the Logos. The Logos brings the phenomenal world into manifestation. In other words, the Logos is the creative cosmic Word that incarnates into the world. However, the things of the world are shadows, different modes of the one true Existence.

\paragraph{Aristotle}


\textbf{Aristotle}, on the other hand, claimed that the idea inheres in the thing, not in some abstract world. The idea and the thing are inseparably connected. This gives reality to the physical world of the senses. This union of idea and substance overcomes the chaos of the manifest world. Like Plato, Aristotle also knew of the primordial idea which is called the Logos. However, it is most real when incarnated into the world of senses.


\flrightit{Posted on 2022-03-03 by Cologero}

